%% 第 1 章：算子与不变子空间
%% 本文件由 tools/md2tex.py 自动生成，请勿直接编辑。

\chapter{算子与不变子空间}


\begin{jdefn}{算子}{r:2:1}
从向量空间 $V$ 映到自身的线性映射：

\[
T:V\to V
\]

称为 $V$ 上的算子。

因为定义域与陪域相同，算子可以反复作用：

\[
v,\quad T(v),\quad T^2(v),\quad T^3(v),\ldots
\]

其中：

\[
T^0=I,
\qquad
T^{k+1}=T\circ T^k.
\]
\end{jdefn}

\begin{jdefn}{不变子空间}{r:2:2}
若 $U\le V$，且：

\[
T(U)\subseteq U,
\]

则称 $U$ 是 $T$ 的不变子空间。

这表示：从 $U$ 中任取一个向量，反复施加 $T$ 后，结果始终不会离开 $U$。
\end{jdefn}

\section{不变子空间为何是分解的基础}

若：

\[
V=U_1\oplus\cdots\oplus U_m,
\]

并且每个 $U_j$ 都对 $T$ 不变，则任意向量可唯一写成：

\[
v=u_1+\cdots+u_m,
\qquad
u_j\in U_j.
\]

由线性性：

\[
T(v)=T(u_1)+\cdots+T(u_m).
\]

而每个 $T(u_j)$ 仍位于对应的 $U_j$ 中。

因此，算子在不同不变子空间上的作用彼此不会混合。我们可以先分别研究 $T|_{U_j}$，再把结论合并。

\begin{jinsight}
算子分解的核心，就是寻找足够多的不变子空间，并判断它们能否构成直和。
\end{jinsight}
